\documentclass[a4paper,12pt]{article}
\setlength{\headheight}{68.803pt}
\addtolength{\topmargin}{-56.803pt}

\usepackage[utf8]{inputenc}
\usepackage[T1]{fontenc}
\usepackage{lmodern}
\usepackage[nopatch=footnote]{microtype}
\usepackage{fvextra}
\DefineVerbatimEnvironment{Verbatim}{Verbatim}{
  breaklines=true,
  breakanywhere=true,
  frame=single,
  fontsize=\small
}

% Impostazione dei margini con geometry
\usepackage[
  left=1in,
  right=1in,
  top=3cm,
  bottom=1in
]{geometry}

\usepackage[
    colorlinks=true,
    linkcolor=black,
    urlcolor=blue,
    citecolor=blue,
    pdfborder={0 0 0},
    pdfauthor={Massimo Mantineo},
    pdftitle={MapReduce Distribuito con Ray.io},
    pdfsubject={Laboratorio di Reti e Sistemi Distribuiti: HANDS-ON 18},
]{hyperref}

\usepackage{graphicx}
\graphicspath{{../../assets/}{../../assets/ho1/}{../../assets/ho2/}{../../assets/ho3/}{../../assets/ho4/}{../../assets/ho5/}{../../assets/ho6/}{../../assets/ho7/}{../../assets/ho8/}{../../assets/ho10/}}
\usepackage{tikz}
\usepackage{listings}
\usepackage{xcolor}
\usepackage{amsmath, amssymb}
\usepackage{fancyhdr}
\usepackage{setspace}
\usepackage{pgfplots}
\pgfplotsset{compat=1.17}
\usepackage{booktabs}
\usepackage[skins, listings]{tcolorbox}
\usepackage{float}

\newtcblisting{terminal}[1][]{
    listing only,
    colback=black!85!white,
    colframe=black!70!white,
    colupper=green!80!white,
    coltitle=white,
    title=#1,
    fonttitle=\bfseries\sffamily,
    listing options={
        language={},
        basicstyle=\ttfamily\small\color{green!80!white},
        backgroundcolor={},
        frame=none,
        breaklines=true,
        showstringspaces=false,
        numbers=none
    },
    enhanced,
    drop shadow,
    arc=3mm
}

\newtcblisting{ccode}[1][]{
    listing only,
    colback=blue!3!white,
    colframe=blue!70!black,
    title=#1,
    fonttitle=\bfseries\sffamily,
    listing options={
        style=cstyle,
        backgroundcolor={},
        frame=none
    },
    enhanced,
    drop shadow,
    arc=3mm
}

\setlength{\footskip}{50pt}

\lstdefinestyle{cstyle}{
    language=C,
    basicstyle=\ttfamily\small,
    numbers=left,
    numberstyle=\tiny,
    stepnumber=1,
    numbersep=5pt,
    backgroundcolor=\color{white},
    showspaces=false,
    showstringspaces=false,
    showtabs=false,
    frame=single,
    tabsize=4,
    captionpos=b,
    breaklines=true,
    breakatwhitespace=true,
    keywordstyle=\color{blue}\bfseries,
    commentstyle=\color{green!50!black},
    stringstyle=\color{red},
}
\lstset{style=cstyle}

% Impostazione di fancyhdr per header e footer
\pagestyle{fancy}
\fancyhf{}

\fancyhead[L]{%
  \parbox[b]{0.45\textwidth}{%
    \small
    Student Name: Massimo Mantineo\\
    Student ID: 541924
    \vspace{20pt}
  }%
}
\fancyhead[C]{%
  \parbox[b]{0.2\textwidth}{%
    \centering
    \normalsize \bfseries
    LabReSiD25\\
    Hands-On 18
    \vspace{20pt}
  }%
}
\fancyhead[R]{%
  \parbox[b]{0.35\textwidth}{%
    \flushright
    \includegraphics[width=3.5cm]{\detokenize{logo_unime_orizontale.png}}
    \vspace{15pt}
  }%
}

\renewcommand{\contentsname}{Indice}

\title{\vspace{-2cm}Report: MapReduce Distribuito con Ray.io\\
\large (Hands-On 18)}
\author{Massimo Mantineo \\ Matricola: 541924}
\date{MESSINA, A.A. 2024/2025}

\begin{document}

% Pagina di copertina
\begin{titlepage}
    \thispagestyle{empty}
    \centering
    \vspace*{1cm}
    \includegraphics[width=6cm]{\detokenize{logo_unime_orizontale.png}}\par\vspace{1cm}
    \vspace{2cm}
    {\Large \textbf{Università degli Studi di Messina}\par}
    {\large Dipartimento MIFT - Corso di Laurea Triennale in Informatica (L-31)\par}
    \vspace*{1cm}
    {\large Laboratorio di Reti e Sistemi Distribuiti\par}
    \vspace{1cm}
    {\Large \textbf{HANDS-ON 18:}\\
    MapReduce Distribuito con Ray.io\par}
    \vspace{2cm}
    {\large Massimo Mantineo \par}
    {\large Matricola: 541924 \par}
    \vfill
    {\large Messina, A.A. 2024/2025 \par}
\end{titlepage}

\clearpage
\pagestyle{fancy}
\tableofcontents
\newpage

\section{Problem Statement}
L'Hands-On 18 affronta lo studio e l'implementazione di un sistema di calcolo distribuito basato sul paradigma \textbf{MapReduce}. A differenza degli esercizi precedenti svolti in C e focalizzati sulla comunicazione inter-processo a basso livello o socket, questo esperimento utilizza \textbf{Python} e il framework \textbf{Ray.io} per gestire la parallelizzazione e la distribuzione dei task su cluster (reali o virtuali).

L'obiettivo specifico è il conteggio delle frequenze dei termini presenti nel testo del celebre romanzo \textit{"Il Signore degli Anelli"}, analizzando come i task vengono suddivisi, eseguiti in parallelo e infine aggregati.

\section{Stato dell'Arte}
Il paradigma MapReduce, reso celebre da Google, si basa su due fasi principali:
\begin{itemize}
    \item \textbf{Map}: I dati di input vengono suddivisi in piccoli pezzi e processati in parallelo da diversi worker.
    \item \textbf{Reduce}: I risultati parziali dei worker vengono aggregati per produrre il risultato finale.
\end{itemize}

Framework moderni come \textbf{Ray} permettono di implementare tali pattern in modo efficiente, offrendo astrazioni semplici per l'esecuzione remota di funzioni e la gestione delle dipendenze, superando i limiti della libreria standard multiprocess di Python.

\section{Metodologia ed Implementazione}
L'architettura del sistema si basa sul modello a tre stadi (Map, Shuffle, Reduce) orchestrato da Ray.

\begin{figure}[H]
    \centering
    \begin{tikzpicture}[
        node distance=1.5cm,
        box/.style={rectangle, draw=blue!70, fill=blue!5, thick, minimum width=2.5cm, minimum height=1cm, align=center, font=\small\sffamily},
        data/.style={draw=orange!70, fill=orange!5, thick, rectangle, rounded corners, font=\tiny\ttfamily},
        arrow/.style={-Stealth, thick, gray!80}
    ]
        % Stages
        \node[data] (input) {Input: lotr.txt};
        \node[box, right=of input] (split) {Splitter\\(Text Chunks)};
        
        \node[box, right=of split, yshift=1cm] (map1) {Map Worker 1\\(Word Counts)};
        \node[box, right=of split, yshift=-1cm] (map2) {Map Worker 2\\(Word Counts)};
        
        \node[box, right=2cm of split] (agg) {Global\\Reducer\\(ray.get)};
        \node[data, right=of agg] (output) {Final Report};

        % Connections
        \draw[arrow] (input) -- (split);
        \draw[arrow] (split) -- (map1);
        \draw[arrow] (split) -- (map2);
        \draw[arrow] (map1) -- (agg);
        \draw[arrow] (map2) -- (agg);
        \draw[arrow] (agg) -- (output);

        % Annotations
        \node[above=0.1cm of map1, font=\tiny\itshape, color=blue!70] {@ray.remote};
        \node[above=0.1cm of map2, font=\tiny\itshape, color=blue!70] {@ray.remote};
    \end{tikzpicture}
    \caption{Workflow del calcolo distribuito MapReduce orchestrato tramite Ray.io.}
    \label{fig:mapreduce_flow}
\end{figure}

L'architettura del sistema prevede:
\begin{enumerate}
    \item \textbf{Splitting}: Il testo di input (\texttt{lotr.txt}) viene suddiviso in $N$ blocchi.
    \item \textbf{Remote Map Tasks}: Ogni blocco è processato da una funzione Ray remota (\texttt{@ray.remote}) che produce un \texttt{Counter} locale.
    \item \textbf{Aggregation}: I risultati vengono scaricati dal cluster tramite \texttt{ray.get()} e fusi nel Counter globale.
\end{enumerate}

Il codice principale è contenuto in \texttt{mapreduce\_lotr.py}.


\subsection{Istruzioni di Esecuzione}
L'esecuzione del sistema MapReduce richiede Python 3:
\begin{terminal}[Run]
$ make setup # Installa dipendenze venv
$ make run   # Avvia l'analisi LOTR
\end{terminal}

\section{Risultati}
L'esecuzione sul testo completo del romanzo ("The Lord of the Rings") ha prodotto i seguenti risultati (Top 10 parole per frequenza):

\begin{terminal}[Output MapReduce Distribuito]
Starting distributed MapReduce on lotr.txt...
Report generated: lotr\_report.txt
Top 10 words:
  the: 11705
  and: 7555
  of: 5083
  to: 3948
  a: 3728
  he: 3038
  in: 2961
  i: 2825
  it: 2780
  that: 2539
\end{terminal}

I risultati mostrano la tipica distribuzione delle lingue naturali, dove articoli e congiunzioni dominano il conteggio. La distribuzione tramite Ray è avvenuta correttamente, inizializzando un cluster locale sulla macchina di test e completando l'analisi di migliaia di parole in pochi secondi.

\section{Conclusioni e Sviluppi Futuri}
L'esperimento dimostra l'efficacia di Ray per l'elaborazione distribuita. Sviluppi futuri potrebbero includere:
\begin{itemize}
    \item Implementazione di fasi di Reduce distribuite (Shuffling).
    \item Utilizzo di dataset massivi caricati da cloud storage (AWS S3/GCS).
    \item Filtraggio di "stop-words" per un'analisi semantica più accurata.
\end{itemize}

\end{document}
